\section{补充结果}
\label{app:results}

\subsection{逐任务结果与回合结局}
\label{app:per_task}
表~\ref{tab:supp_per_task} 把表~\mainref{tab:main} 中 \bench 的留出结果按任务拆分到每个方法，并给出 \ours 在每个任务上相对最强其他方法的配对区间。表~\ref{tab:supp_failures} 报告回合如何结束。一个方法若以一种失败换另一种失败，例如以更少的碰撞换来更多的超时，只能在这里看出来。所有失败都留在成功率的分母中。

% Wrap each Delta cell as \dcell{value}{text} once results arrive (shading rule in the caption).
% Transcribed from results/realtime_c/ ACT/DP per-task CSV and reports (INFER0); other methods pending.
\begin{table*}[t]
\centering\footnotesize
\caption{\tablelang{\textbf{Per-task held-out success} (\%) on \bench for every method of Table~\mainref{tab:main}, mean of
three training seeds for our variants and one for baselines, real-time protocol. DP: Diffusion Policy. DVLA: DynamicVLA. DrVLA: DreamVLA. OWAM: OpenWAM-$\alpha$.
DWAM: DynamicWAM. +: with our SE(3) history tokens. $\sigma$: \ours with
nominal synthetic correlated pose noise (Appendix~\ref{app:pose_noise}). $\Delta$: \ours minus the strongest other method on that task, with its 95\% paired interval. The last two rows give the position error at first contact $e_p$ (mm), averaged over the 13 tasks with a contact, and the number $n$ of episodes with a contact (Appendix~\ref{app:metrics}); --: no 13-task mean. Task codes are
listed in Table~\ref{tab:roles}. Shading of $\Delta$: blue for higher SR, orange for lower, saturating at 15 points.}{\textbf{\bench 各任务的留出成功率}（\%），涵盖表~\mainref{tab:main} 中的全部方法，本文各变体取三个训练种子的平均、基线训练一次，实时协议。
DP：Diffusion Policy。DVLA：DynamicVLA。DrVLA：DreamVLA。OWAM：OpenWAM-$\alpha$。DWAM：DynamicWAM。+：加入本文 SE(3) 历史 token。
$\sigma$：使用标称合成相关位姿噪声的 \ours（附录~\ref{app:pose_noise}）。$\Delta$：该任务上 \ours 相对最强其他方法之差及其 95\% 配对区间。最后两行为首次接触时的位置误差 $e_p$（mm，在有接触的 13 个任务上平均）与发生接触的回合数 $n$（附录~\ref{app:metrics}）；--：无 13 任务均值。
任务代号见表~\ref{tab:roles}。$\Delta$ 的底色：蓝色表示 SR 更高，橙色表示更低，在 15 个百分点处饱和。}}
\label{tab:supp_per_task}
\setlength{\tabcolsep}{2.4pt}
\begin{tabular}{@{}lccccccccccccc@{}}
\toprule
\tablelang{Task}{任务} & ACT & DP & $\pi_{0.5}$ & $\pi_{0.5}$+ & DVLA & PUMA & DrVLA & OWAM & DWAM & DWAM+ & $\sigma$ & \textbf{\tablelang{Ours}{本文}} & $\Delta$ [95\% CI] \\
\midrule
\tablelang{Harmonic grasp}{谐波抓取} & 48.0 & 0.0 & \resultplaceholder & \resultplaceholder & \resultplaceholder & \resultplaceholder & \resultplaceholder & \resultplaceholder & \resultplaceholder & \resultplaceholder & \resultplaceholder & \resultplaceholder & \resultplaceholder[{+00.0 [00.0, 00.0]}] \\
\tablelang{Conveyor grasp}{传送带抓取} & 46.0 & 0.0 & \resultplaceholder & \resultplaceholder & \resultplaceholder & \resultplaceholder & \resultplaceholder & \resultplaceholder & \resultplaceholder & \resultplaceholder & \resultplaceholder & \resultplaceholder & \resultplaceholder[{+00.0 [00.0, 00.0]}] \\
\tablelang{Translating grasp}{平移抓取} & 18.0 & 0.0 & \resultplaceholder & \resultplaceholder & \resultplaceholder & \resultplaceholder & \resultplaceholder & \resultplaceholder & \resultplaceholder & \resultplaceholder & \resultplaceholder & \resultplaceholder & \resultplaceholder[{+00.0 [00.0, 00.0]}] \\
\tablelang{Rotating grasp}{旋转抓取} & 21.0 & 0.0 & \resultplaceholder & \resultplaceholder & \resultplaceholder & \resultplaceholder & \resultplaceholder & \resultplaceholder & \resultplaceholder & \resultplaceholder & \resultplaceholder & \resultplaceholder & \resultplaceholder[{+00.0 [00.0, 00.0]}] \\
\tablelang{Occluded grasp}{遮挡抓取} & 32.0 & 0.0 & \resultplaceholder & \resultplaceholder & \resultplaceholder & \resultplaceholder & \resultplaceholder & \resultplaceholder & \resultplaceholder & \resultplaceholder & \resultplaceholder & \resultplaceholder & \resultplaceholder[{+00.0 [00.0, 00.0]}] \\
\midrule
\tablelang{Moving receptacle}{移动容器} & 26.0 & 0.0 & \resultplaceholder & \resultplaceholder & \resultplaceholder & \resultplaceholder & \resultplaceholder & \resultplaceholder & \resultplaceholder & \resultplaceholder & \resultplaceholder & \resultplaceholder & \resultplaceholder[{+00.0 [00.0, 00.0]}] \\
\tablelang{Rotating receptacle}{旋转容器} & 0.0 & 0.0 & \resultplaceholder & \resultplaceholder & \resultplaceholder & \resultplaceholder & \resultplaceholder & \resultplaceholder & \resultplaceholder & \resultplaceholder & \resultplaceholder & \resultplaceholder & \resultplaceholder[{+00.0 [00.0, 00.0]}] \\
\midrule
\tablelang{Wire passage}{导线穿越} & 99.0 & 0.0 & \resultplaceholder & \resultplaceholder & \resultplaceholder & \resultplaceholder & \resultplaceholder & \resultplaceholder & \resultplaceholder & \resultplaceholder & \resultplaceholder & \resultplaceholder & \resultplaceholder[{+00.0 [00.0, 00.0]}] \\
\tablelang{Wire passage, full}{导线穿越（完整）} & 0.0 & 0.0 & \resultplaceholder & \resultplaceholder & \resultplaceholder & \resultplaceholder & \resultplaceholder & \resultplaceholder & \resultplaceholder & \resultplaceholder & \resultplaceholder & \resultplaceholder & \resultplaceholder[{+00.0 [00.0, 00.0]}] \\
\midrule
\tablelang{Ball capture}{小球捕获} & 36.0 & 0.0 & \resultplaceholder & \resultplaceholder & \resultplaceholder & \resultplaceholder & \resultplaceholder & \resultplaceholder & \resultplaceholder & \resultplaceholder & \resultplaceholder & \resultplaceholder & \resultplaceholder[{+00.0 [00.0, 00.0]}] \\
\tablelang{Ball capture, table}{小球捕获（桌面）} & 71.0 & 0.0 & \resultplaceholder & \resultplaceholder & \resultplaceholder & \resultplaceholder & \resultplaceholder & \resultplaceholder & \resultplaceholder & \resultplaceholder & \resultplaceholder & \resultplaceholder & \resultplaceholder[{+00.0 [00.0, 00.0]}] \\
\tablelang{Slider capture}{滑块捕获} & 57.0 & 0.0 & \resultplaceholder & \resultplaceholder & \resultplaceholder & \resultplaceholder & \resultplaceholder & \resultplaceholder & \resultplaceholder & \resultplaceholder & \resultplaceholder & \resultplaceholder & \resultplaceholder[{+00.0 [00.0, 00.0]}] \\
\tablelang{Can capture}{罐体捕获} & 90.0 & 0.0 & \resultplaceholder & \resultplaceholder & \resultplaceholder & \resultplaceholder & \resultplaceholder & \resultplaceholder & \resultplaceholder & \resultplaceholder & \resultplaceholder & \resultplaceholder & \resultplaceholder[{+00.0 [00.0, 00.0]}] \\
\midrule
\tablelang{Post-impact grasp}{撞击后抓取} & 44.0 & 0.0 & \resultplaceholder & \resultplaceholder & \resultplaceholder & \resultplaceholder & \resultplaceholder & \resultplaceholder & \resultplaceholder & \resultplaceholder & \resultplaceholder & \resultplaceholder & \resultplaceholder[{+00.0 [00.0, 00.0]}] \\
\midrule
\tablelang{Task-balanced}{任务等权} & 42.0 & 0.0 & \resultplaceholder & \resultplaceholder & \resultplaceholder & \resultplaceholder & \resultplaceholder & \resultplaceholder & \resultplaceholder & \resultplaceholder & \resultplaceholder & \resultplaceholder & \resultplaceholder[{+00.0 [00.0, 00.0]}] \\
\midrule
$e_p$ (mm) & 32.2 & -- & \resultplaceholder & \resultplaceholder & \resultplaceholder & \resultplaceholder & \resultplaceholder & \resultplaceholder & \resultplaceholder & \resultplaceholder & \resultplaceholder & \resultplaceholder & -- \\
$n$ & 1109 & 6 & \resultplaceholder & \resultplaceholder & \resultplaceholder & \resultplaceholder & \resultplaceholder & \resultplaceholder & \resultplaceholder & \resultplaceholder & \resultplaceholder & \resultplaceholder & -- \\
\bottomrule
\end{tabular}
\end{table*}

% Transcribed from results/realtime_c/act-report.json (INFER0); other methods pending.
\begin{table}[t]
\centering\footnotesize
\caption{\tablelang{\textbf{Outcome categories} on held-out \bench scenes (\% of episodes, seed-averaged, real-time protocol).
Categories follow Appendix~\ref{app:tasks}, and rows sum to 100 before rounding. The noisy-pose row reuses the nominal synthetic correlated-noise episodes from the robustness study.}{\textbf{结局类别}：\bench 留出场景（各类 episode 占比，\%，种子平均，实时协议）。
类别定义见附录~\ref{app:tasks}；各行在四舍五入前合计为 100。含噪位姿行复用鲁棒性实验中标称合成时序相关噪声的轨迹。}}
\label{tab:supp_failures}
\setlength{\tabcolsep}{2.2pt}
\begin{tabular}{@{}lcccccc@{}}
\toprule
\tablelang{Method}{方法} & \tablelang{Success}{成功} & \tablelang{Safety}{安全} & \tablelang{Task}{任务违规} & \tablelang{Bound.}{边界} & \tablelang{Timeout}{超时} & \tablelang{Inval.}{无效} \\
\midrule
ACT & 42.0 & 20.5 & 3.9 & 16.1 & 17.6 & 0.0 \\
Diffusion Policy & \resultplaceholder & \resultplaceholder & \resultplaceholder & \resultplaceholder & \resultplaceholder & \resultplaceholder \\
$\pi_{0.5}$ & \resultplaceholder & \resultplaceholder & \resultplaceholder & \resultplaceholder & \resultplaceholder & \resultplaceholder \\
$\pi_{0.5}$ + SE(3) & \resultplaceholder & \resultplaceholder & \resultplaceholder & \resultplaceholder & \resultplaceholder & \resultplaceholder \\
DynamicVLA & \resultplaceholder & \resultplaceholder & \resultplaceholder & \resultplaceholder & \resultplaceholder & \resultplaceholder \\
PUMA & \resultplaceholder & \resultplaceholder & \resultplaceholder & \resultplaceholder & \resultplaceholder & \resultplaceholder \\
DreamVLA & \resultplaceholder & \resultplaceholder & \resultplaceholder & \resultplaceholder & \resultplaceholder & \resultplaceholder \\
OpenWAM-$\alpha$ & \resultplaceholder & \resultplaceholder & \resultplaceholder & \resultplaceholder & \resultplaceholder & \resultplaceholder \\
DynamicWAM & \resultplaceholder & \resultplaceholder & \resultplaceholder & \resultplaceholder & \resultplaceholder & \resultplaceholder \\
DynamicWAM + SE(3) & \resultplaceholder & \resultplaceholder & \resultplaceholder & \resultplaceholder & \resultplaceholder & \resultplaceholder \\
\midrule
\tablelang{\ours, CV input}{\ours，匀速外推输入} & \resultplaceholder & \resultplaceholder & \resultplaceholder & \resultplaceholder & \resultplaceholder & \resultplaceholder \\
\tablelang{\ours, noisy poses}{\ours，含噪位姿} & \resultplaceholder & \resultplaceholder & \resultplaceholder & \resultplaceholder & \resultplaceholder & \resultplaceholder \\
\ours & \resultplaceholder & \resultplaceholder & \resultplaceholder & \resultplaceholder & \resultplaceholder & \resultplaceholder \\
\bottomrule
\end{tabular}
\end{table}


\subsection{协议与延迟}
\label{app:protocols}
表~\ref{tab:supp_protocols} 在另外三种协议下评估每个方法。暂停物理仿真去除了计算的代价。等延迟协议向每个方法计入 \ours 的 p95 延迟，从而把方法的速度与其利用时间的方式区分开。原执行器协议检验我们的共享执行器是否使某个方法处于不利地位。在暂停物理仿真时，\ours 与 \fillin{DreamVLA} 的差距为 \fillin{$\pm$X.X} 个百分点，在等延迟下为 \fillin{$\pm$X.X} 个百分点。原执行器使各基线的 SR 至多变化 \fillin{X.X} 个百分点。表~\ref{tab:supp_latency} 给出扫描中每个附加延迟下的成功率，其中 \ours 与最强基线的两行即图~\mainref{fig:teaser}（右）的数据。

% Transcribed from results/realtime_c/formal64-20261101.json (Research 2d190ac55); final Ours and unmeasured protocols remain placeholders.
\begin{table}[H]
\centering\footnotesize
\caption{\tablelang{\textbf{Protocol controls} on held-out \bench scenes (SR, \%) under the protocols of
Appendix~\ref{app:sweeps}; real time is the protocol of Table~\mainref{tab:main}. Late: queries in real time whose
time-to-action exceeds the 50\,ms control period (\%). Age: p95 action age in real time (ms); --: no command executed. Blue/bold and light blue rank the reported values only; pending entries are excluded.}{\textbf{评测协议对照}：\bench 留出场景（SR，\%），协议定义见附录~\ref{app:sweeps}；实时即表~\mainref{tab:main} 的协议。超时：实时协议下出动作时间超过 50\,ms 控制周期的查询比例（\%）。年龄：实时协议下动作年龄的 p95（ms）；--：没有执行任何指令。 粗体加蓝底与浅蓝底仅对已有结果排名，待测项不参与。}}
\label{tab:supp_protocols}
\setlength{\tabcolsep}{3pt}
\begin{tabular}{@{}lccccc@{}}
\toprule
 & \tablelang{Equal}{等} & \tablelang{Original}{原执行} & \tablelang{Real}{实时} & \tablelang{Late}{超时} & \tablelang{Age}{年龄} \\
\tablelang{Method}{方法} & \tablelang{latency}{延迟} & \tablelang{executor}{执行器} & \tablelang{time}{} & (\%) & p95 \\
\midrule
ACT & \resultplaceholder & \resultplaceholder & 41.6 & 0.0 & 50 \\
Diffusion Policy & \resultplaceholder & \resultplaceholder & 30.7 & 100.0 & 150 \\
$\pi_{0.5}$ & \resultplaceholder & \resultplaceholder & \secondcell{53.8} & 100.0 & 150 \\
$\pi_{0.5}$ + SE(3) & \resultplaceholder & \resultplaceholder & 52.8 & 100.0 & 150 \\
DynamicVLA & \resultplaceholder & \resultplaceholder & 50.9 & 100.0 & 150 \\
PUMA & \resultplaceholder & \resultplaceholder & \bestcell{59.6} & 100.0 & 150 \\
OpenWAM-$\alpha$ & \resultplaceholder & \resultplaceholder & 46.1 & 100.0 & 350 \\
OpenWAM-$\alpha$ (off) & \resultplaceholder & \resultplaceholder & 15.8 & 100.0 & 950 \\
DynamicWAM & \resultplaceholder & \resultplaceholder & 40.5 & 100.0 & 250 \\
DynamicWAM + SE(3) & \resultplaceholder & \resultplaceholder & 36.8 & 100.0 & 350 \\
\midrule
\ours & \resultplaceholder & -- & \resultplaceholder & \resultplaceholder & \resultplaceholder \\
\bottomrule
\end{tabular}
\end{table}

% Transcribed from results/realtime_c/latency-five-point-20261104.json: eleven configurations, 55 accepted points; one training seed each.
\begin{table}[H]
\centering\footnotesize
\caption{\tablelang{\textbf{Success versus added latency} on 1{,}400 held-out \bench scenes per point (SR in \%, real-time protocol). The 0 and
150\,ms columns equal the Native and +150\,ms columns of Table~\mainref{tab:main}. Blue/bold and light blue mark the best and second-best values per column.}{\textbf{成功率随附加延迟的变化}：每点使用 1{,}400 个 \bench 留出场景（SR，\%；实时协议）。0 与 150\,ms 两列分别等于表~\mainref{tab:main}
的原生与 +150\,ms 列。 粗体加蓝底与浅蓝底分别标记各列最佳与次佳。}}
\label{tab:supp_latency}
\setlength{\tabcolsep}{3pt}
\begin{tabular}{@{}lccccc@{}}
\toprule
 & \multicolumn{5}{c}{\tablelang{Added latency (ms)}{附加延迟（ms）}} \\
\cmidrule(lr){2-6}
\tablelang{Method}{方法} & 0 & 50 & 100 & 150 & 200 \\
\midrule
ACT & 41.6 & 43.5 & \secondcell{36.5} & 24.3 & 16.5 \\
Diffusion Policy & 30.7 & 26.1 & 21.0 & 15.9 & 14.3 \\
$\pi_{0.5}$ & 53.8 & \secondcell{47.2} & \bestcell{39.6} & 27.8 & 13.6 \\
$\pi_{0.5}$ + SE(3) & 52.8 & 39.8 & 32.4 & \secondcell{28.2} & 15.1 \\
DynamicVLA & 50.9 & 46.4 & 30.6 & 23.4 & 17.8 \\
PUMA & \secondcell{59.6} & \bestcell{53.2} & 36.4 & \bestcell{31.1} & \bestcell{27.1} \\
OpenWAM-$\alpha$ & 46.1 & 36.1 & 27.9 & 24.7 & \secondcell{22.1} \\
OpenWAM-$\alpha$ (off) & 15.8 & 13.7 & 14.2 & 10.4 & 9.5 \\
DynamicWAM & 40.5 & 27.6 & 25.0 & 20.6 & 20.9 \\
DynamicWAM + SE(3) & 36.8 & 30.6 & 23.9 & 21.9 & 21.3 \\
\midrule
\ours & \bestcell{64.3} & 43.4 & 25.5 & 9.9 & 6.6 \\
\bottomrule
\end{tabular}
\end{table}


\subsection{负面与无差异结果}
\label{app:negative}

\paragraph{筛选结果。}
历史、预测目标、编码器与耦合比较采用附录~\ref{app:design} 的筛选配方。每项结果必须同时注明其原始划分与种子数。Q3 的耦合效应在种子之间变号，不能证明最终规模下训练监督与在线预测反馈等效。新的四对照研究在表~\mainref{tab:design} 中报告带符号差值与配对区间。\fillin{按实际规模概括已完成的零效应与负面效应。}

表~\ref{tab:supp_remember} 与~\ref{tab:supp_predict} 按原始规模报告已完成的历史与预测筛选结果。表~\ref{tab:supp_interventions} 报告冻结筛选权重上的静止历史干预；它不是最终规模的训练消融。

\paragraph{未解决的任务。}放入旋转料箱的 SE(3) 放置任务（t2\_se3）对所有方法都失败（表~\ref{tab:supp_per_task}）。\fillin{策略能到达料箱，但在它转开之前达不到其角度容差。}\ours 的任何变体都没有解决完整的导线穿越任务（t3\_full），该任务在穿越之外还要求抓取与释放棒，而 \fillin{OpenWAM-$\alpha$} 在其中 \fillin{XX}\% 的回合中成功。在以已抓住棒开始的穿越任务（t3\_phase）上，\ours 的成功率在不同训练种子之间从 \fillin{0} 到 \fillin{71}\% 不等。

\paragraph{其他失败。}图~\ref{fig:budget_trends}a 中的曲线聚合在捕获任务上最弱（Spearman $\rho=\fillin{X.XX}$），因为这些任务的成功还取决于夹爪何时合拢。三个合成噪声条件使 SR 变化\fillin{报告相对无噪声参考的三个带符号差值}；这些结果与真实跟踪器标定分开。\ours 的其余失败大多是\fillin{例如首次抓取失手后的超时}（表~\ref{tab:supp_failures}）。

\subsection{执行时刻的预测}
\label{app:qualitative}
\begin{figure*}[t]
\centering
\figplaceholder[2.1in]{\textwidth}{图 S5——与实际发生对齐的预测}{%
\textbf{第 1--4 行：}每种运动类型（传送带、旋转目标、滚动捕获、撞击后）预测误差居中的留出回合；\textbf{第 5 行：}一次失败。
\textbf{列：}$\tau_t$ 处的查询帧；首条指令执行时刻 $\tapp$ 处的帧；基座坐标系中的 3D 插图。\textbf{叠加：}\textcolor{obscol}{蓝色}过去的 SE(3) 位姿；\textcolor{predcol}{橙色}$\tau_t+\delta_h$ 处预测的目标坐标轴与稀疏流箭头（读自训练期预测头）；\textcolor{gtcol}{灰色}同一物理时刻的真值位姿；\textcolor{actcol}{紫色}发出的夹爪轨迹；首次接触处的标记。每行下方的条带画出学习的 SE(3)、匀速与匀加速预测的平移与旋转误差随时域的变化。用\textcolor{failcol}{红色}标注第 5 行的失败原因。}
\caption{训练期预测头在留出回合上的预测，在物理时间上与真值和发出的动作对齐。第 1--4 行为每种运动类型预测误差居中的回合，第 5 行为一次失败。}
\label{fig:qualitative}
\end{figure*}

图~\ref{fig:qualitative} 把查询、预测、真值、发出的动作与首次接触按物理时间戳对齐。这些示例展示预测误差出现的位置，例如撞击的瞬间或滚动物体离开斜坡时，并同时给出接触时刻与位姿误差。这些对齐的诊断不证明部署的动作流读取了解码预测，也不证明某个预测误差导致了某次动作误差。

\subsection{补充视频}
\label{app:video}
补充视频展示 \ours 与最强基线在全部六个真机任务上的试验（叠加跟踪位姿）、每个任务族的一组 \bench 回合（叠加训练期预测头的预测），以及第~\mainref{sec:real}~节讨论的失败案例。
